
\begin{proposition}[Lema de Abhyankar relativo]
\label{XIII.5.5}
\index{Abhyankar relativo (lema de)|hyperpage}%
Sean $X$ un $S$-esquema y $D$ un divisor de cruzamientos normales
relativamente a~$S$, como en~\eqref{XIII.5.4}. Sean
\ifthenelse{\boolean{orig}}
{$Y=X-\Supp D$}
{$Y=\Supp D$},
$U=X-Y$, $x$ un punto de~$Y$, $X_{1}$ la localización estricta de~$X$
en un punto geométrico situado sobre~$x$, $U_{1}=U\times_{X}X_{1}$ y
$V_{1}$ un recubrimiento étale de~$U_{1}$. Se supone que, para todo
punto maximal $s$ de~$S$, $V_{1\overline{s}}$ está moderadamente
ramificado sobre $X_{1\overline{s}}$ relativamente a~$\overline{s}$.
Entonces pueden encontrarse enteros $n_{i}$, primos con la
característica~$p$ de~$\kres(x)$, con
$i\in I(x)$, tales que, si se pone
$$
X'_{1}=X_{1}[T_{i}]_{i\in I(x)}\big/(T_{i}^{n_{i}}-f_{i})\quoi,
$$
$U'_{1}=U_{1}\times_{X_{1}}X'_{1}$, etc., el recubrimiento étale
$V'_{1}$ de~$U'_{1}$ se prolonga, de manera única salvo isomorfismo
único, a un recubrimiento étale de~$X'_{1}$. En particular,
\marginpar{432}
$V_{1}$ está moderadamente ramificado sobre~$X_{1}$ relativamente a~$S$.
\end{proposition}

Puede suponerse que $S$ es local noetheriano, con punto cerrado~$f(x)$.
Para cada punto maximal $s$ de~$S$ y cada $i\in I(x)$, sea $x_{i}$ el
punto genérico del cerrado $V(f_{i})$ de la fibra $X_{1\overline{s}}$.
El anillo local $(\cal{O}_{X_{1},x_{i}})_{\red}$ es un anillo de
valoración discreta con cuerpo de fracciones $K_{i}$, y se tiene
\ifthenelse{\boolean{orig}}
{$V_{1}|K_{i}=\Spec(\prod_{j\in J(x_{i})}L_{j})$,}
{$V_{1}|K_{i}=\Spec(\prod_{j\in J(x_{i})}L_{j})$,}
donde $L_{j}$ es una extensión finita separable de
\ifthenelse{\boolean{orig}}
{$K$}
{$K_{i}$};
se denota por $n_{j}$ el orden del grupo de inercia de una extensión
galoisiana generada por~$L_{j}$, y por $n_{i}$ el m.c.m. de los $n_{j}$
cuando $s$ recorre los puntos maximales de~$S$ y
\ifthenelse{\boolean{orig}}
{$j\in J(x_{i})$.}
{$j\in J(x_{i})$.}

Elegidos así los $n_{i}$, demostraremos que $V'_{1}$ se prolonga de
manera única a un recubrimiento étale de~$X'_{1}$. La unicidad se
deduce de que, como
\ifthenelse{\boolean{orig}}
{$X'$ es liso sobre~$S$ en los puntos de~$Y$}
{$X'$ es liso sobre~$S$ en los puntos de~$Y'$},
se tiene $\prof\et_{Y'_{1}}(X'_{1})\geqslant 2$ (SGA~4 XVI~3.2 o
SGA~2 XIV~1.19). Sean $x'_{1}$ un punto de~$Y'_{1}$,
$\overline{x}'_{1}$ un punto geométrico situado sobre~$x'_{1}$, y
denotemos por $\overline{X'_{1}}$ la localización estricta de~$X'_{1}$
en~$\overline{x}'_{1}$, $\overline{U'_{1}}=U'_{1(\overline{X'_{1}})}$,
etc. Por descenso, habida cuenta de la unicidad, basta demostrar que
$\overline{V}'_{1}$ se prolonga a~$\overline{X'_{1}}$. Además, puede
uno limitarse a tomar como $x'_{1}$ los puntos maximales de~$Y'_{1}$;
en efecto, se tendrá entonces una prolongación de~$V'_{1}$ sobre un
abierto $W'_{1}$ de~$X'_{1}$ que contiene los puntos maximales
de~$Y'_{1}$. Ahora bien, si $Z'_{1}=X'_{1}-W'_{1}$, se tiene
$\codim(Z'_{1s},X'_{1s})\geqslant 2$ cuando $s$ es un punto maximal
de~$S$, y $\codim(Z'_{1s},X'_{1s})\geqslant 1$,
$\prof\et_{s}(S)\geqslant 1$ cuando $s$ es un punto de~$S$ que no es
maximal; el hecho de que $V'_{1}$ se prolongue a todo~$X'_{1}$ se deduce
entonces de (SGA~2 XIV~1.20). Pero, en un punto geométrico
$\overline{x}'_{1}$ situado sobre un punto maximal de~$Y'_{1}$,
\ifthenelse{\boolean{orig}}
{$X'_{1\red}$}
{$\overline{X'_{1}}_{\red}$}
es el espectro de un anillo de valoración discreta, y el hecho de que
$\overline{V'_{1}}$ se prolongue a~$\overline{X'_{1}}$ se deduce entonces
de~\Ref{X.3.6}.

Demostremos que los $n_{i}$ son primos con~$p$. En efecto, si no fuera
así, habría un índice $i_{0}\in I(x)$ tal que $p$ divide a~$n_{i_{0}}$.
Tras reemplazar, si es preciso, $X$ por
\ifthenelse{\boolean{orig}}
{$X_{1}[T_{i_{0}},T_{i}]_{i\in I(x)}\big/
(T_{i}^{n_{i_{0}}/p}-f_{i_{0}},T_{i}^{n_{i}}-f_{i})$,}
{$X_{1}[T_{i}]_{i\in I(x)}\big/
(T_{i_{0}}^{n_{i_{0}}/p}-f_{i_{0}},
(T_{i}^{n_{i}}-f_{i})_{i\in I(x),\,i\ne i_{0}})$,}
se reduce al caso en que
\ifthenelse{\boolean{orig}}
{$X'_{1}=X_{1}[T]/T^{p}-f_{i_{0}}$}
{$X'_{1}=X_{1}[T]/(T^{p}-f_{i_{0}})$}.
Por lo que
precede, el recubrimiento
\marginpar{433}
étale $V'_{1}$ de~$U'_{1}$ se prolonga a un recubrimiento étale $E'_{1}$
de~$X'_{1}$. Sea $\eta$ el punto cerrado de~$S$; como el morfismo
$X'_{1\eta}\to X_{1\eta}$ es radicial, $V_{1\eta}$ se prolonga a un
recubrimiento étale $E_{1\eta}$ de~$X_{1\eta}$. Se deduce entonces,
como en~\Ref{XIII.5.2}, que $E'_{1}$ está provisto de un dato de descenso
relativamente al morfismo $X'_{1}\to X_{1}$, que prolonga el dato de
descenso natural que se tiene sobre~$E'_{1}|U'_{1}$. De aquí resulta que
$V_{1}$ se prolonga a~$X_{1}$; pero esto implica $n_{i_{0}}=1$, en
contra de la hipótesis $n_{i_{0}}=p$.
